% SWITCH-CUT-MENGER3
\\section*{SWITCH-CUT-MENGER3: topology-independent passive contact lower bound}

\\textbf{STATUS: PROVED/EXACT in the declared passive monotone contact model plus COMPUTATIONAL exact finite evaluation. P versus NP remains OPEN.}

Let \\(D(n,8)\\) be the exact-weight-eight clause-presence classifier and write
\\[ M=8\\binom n3. \\]
A passive monotone contact scheme is an undirected two-terminal multigraph whose controlled
contacts are labelled by clause-presence variables and conduct exactly when the corresponding
clause bit is one. Labels may repeat. Fixed always-on wiring is free and may be contracted.
Complement contacts are outside this theorem.

The exact number of UNSAT promised inputs is
\\[
U_8(n)=\\binom n3+268\\binom n4+3030\\binom n5+7800\\binom n6+5040\\binom n7.
\\]
Define \\(\\tau(n)\\) as the least integer \\(d\\) satisfying
\\[ \\binom{M-d}{8}\\le U_8(n). \\]

\\paragraph{Cut-label lemma.}
If an \\(s\\)-\\(t\\) cut contains contacts carrying only \\(d\\) distinct clause labels,
then every legal weight-eight input chosen from the other \\(M-d\\) labels disables that cut.
Hence all \\(\\binom{M-d}{8}\\) such inputs would have to be UNSAT. Therefore every cut
contains at least \\(\\tau(n)\\) distinct labels and at least that many contact edges.

By edge-Menger there are at least \\(\\tau(n)\\) pairwise edge-disjoint \\(s\\)-\\(t\\) paths.
Every clause label belongs to an UNSAT exact-eight formula, namely the eight sign patterns on
its underlying variable triple. Consequently no accepting path can contain only one controlled
contact. After contracting fixed wiring, every path contains at least two controlled contacts.
Thus
\\[ \\boxed{C_{\\rm passive}(n,8)\\ge2\\tau(n).} \\]

Exact evaluation gives
\\[42\\le C_{\\rm passive}(4,8)\\le48,\\qquad
132\\le C_{\\rm passive}(5,8)\\le144.\\]
The upper bounds are the endpoint-routed constructions from SWITCH-RECTIFIER2.

For \\(n=4\\), \\(\\binom{16}{8}=12870>272\\); for \\(n=5\\) the analogous inequality is immediate.
For \\(n\\ge6\\),
\\[
\\frac M2=4\\binom n3\\ge n^3/3,\\qquad
\\binom{M/2}{8}\\ge(n^3/24)^8,
\\]
while
\\[ U_8(n)\\le(581/12)n^7<49n^7. \\]
Since \\(6^{17}>49\\cdot24^8\\), it follows that \\(\\tau(n)>M/2\\) for every \\(n\\ge4\\).
Hence \\(C_{\\rm passive}>M\\): a one-occurrence/read-once passive switch network is
impossible even on the exact-weight-eight promise.

Letting \\(r=M-\\tau\\), the inequalities \\(\\binom r8\\le U_8(n)\\) and
\\(\\binom r8\\ge(r/8)^8\\) give \\(r<13.013n^{7/8}\\). Therefore
\\[
\\boxed{C_{\\rm passive}(n,8)\\ge2M-26.026n^{7/8}.}
\\]
Together with the explicit upper bound \\(2M-16\\), passive contact complexity is asymptotically
\\(2M-o(M)\\) in this fixed-eight family.

A separate exhaustive control verifies that every one of the 96 \\(D(4,8)\\) conflict pairs
has a legal weight-eight SAT input in which it is the unique conflict pair. Thus a simple bank
of two-contact paths cannot sparsify the conflict graph: all 96 pair paths are individually
essential in that subclass.

\\paragraph{Literature context.}
Read-once switch networks are a standard synthesis model. Jukna's entropy and null-chain work
shows that general contact-scheme compression comes from nontrivial topological merging, while
Menger's theorem supplies the cut-to-edge-disjoint-path bridge used here.

\\paragraph{Limitations.}
The theorem is restricted to passive monotone contacts. Normally-closed/complement contacts,
active or sequential relays, general Boolean circuits, branching programs, and RAM algorithms
are stronger models. No unrestricted SAT-time lower bound follows.

\\paragraph{Next target.}
For general \\(D(n,m)\\), with \\(N=8\\binom n3\\) possible clause labels and \\(U(n,m)\\)
UNSAT promised inputs, the same argument yields \\(\\binom{N-d}{m}\\le U(n,m)\\) for every
cut using \\(d\\) distinct labels. Evaluate the induced \\(\\tau(n,m)\\) for \\(m=\\alpha n\\),
especially near the random-3-SAT transition.

\\begin{thebibliography}{9}
\\bibitem{YoshidaIkedaAsada} H.~Yoshida, M.~Ikeda, and K.~Asada,
\\emph{Synthesis of Read-Once Switch Network}, IEICE Society Conference, 2006.
\\bibitem{Jukna88} S.~P.~Jukna, \\emph{Entropy of Contact Circuits and Lower Bounds on Their Complexity},
Theoretical Computer Science 57 (1988), 113--129.
\\bibitem{Jukna89} S.~P.~Jukna, \\emph{The Effect of Null-Chains on the Complexity of Contact Schemes},
FCT 1989, 246--256.
\\bibitem{Menger27} K.~Menger, \\emph{Zur allgemeinen Kurventheorie}, Fundamenta Mathematicae 10 (1927), 96--115.
\\end{thebibliography}

% END_OF_NOTE::SWITCH-CUT-MENGER3